\documentclass[../main]{subfiles}
\begin{document}

\chapter{Primera Parte}
\section{Dibujo Técnico}
\subsection{Introducción(Lectura)}
\subsubsection*{¿Qué es el dibujo técnico?}

El dibujo técnico es un sistema de representación gráfica de diversos tipos de objetos, con el propósito de proporcionar información suficiente para facilitar su análisis, ayudar a elaborar su diseño y posibilitar su futura construcción y mantenimiento. Suele realizarse con el auxilio de medios informatizados o, directamente, sobre el papel u otros soportes planos.

\begin{figure}[ht]
    \centering
    \includegraphics{images/250px-Euro_Construction.svg.png}
    \caption{Dibujo técnico del símbolo monetario del euro.}
\end{figure}

La representación gráfica se basa en la geometría descriptiva y utiliza las proyecciones ortogonales para dibujar las distintas vistas de un objeto.

Los objetos, piezas, máquinas, edificios, planos urbanos, entre otros , se suelen representar en planta (vista superior, vista de techo, planta de piso, cubierta, entre otros ), alzado (vista frontal o anterior) y lateral (acotaciones); son necesarias un mínimo de dos proyecciones (vistas del objeto) para aportar información útil del objeto, dependiendo esto de la complejidad del mismo. Las vistas mencionadas de acuerdo al sistema ortogonal se llaman fundamentales por pertenecer al triedro fundamental, este triedro lo conforman el plano anterior, superior y lateral.

Con el objetivo de unificar el lenguaje del dibujo técnico se establecieron normativas aprobadas internacionalmente, pero cada país tiene su organismo nacional de normalización para el estudio y aprobación de las diferentes Normas (en Argentina es IRAM la única organización que realiza esta tarea). Aplicando estas normativas cualquier plano podrá ser interpretado por cualquier profesional del área correspondiente.

\subsubsection*{Tipos de Dibujo Técnico.}
Con el desarrollo industrial y los avances tecnológicos el dibujo ha aumentado su campo de acción, los principales son:
\begin{itemize}
\item Dibujo arquitectónico

Vistas estándar usadas en el dibujo arquitectónico.
El dibujo arquitectónico abarca una gama de representaciones gráficas con las cuales se realizan los planos para la construcción de edificios, casas, quintas, autopistas, iglesias, fábricas y puentes entre otros. Se dibuja el proyecto con instrumentos precisos, con sus respectivos detalles, ajuste y correcciones, donde aparecen los planos de planta, fachadas, secciones, perspectivas, fundaciones, columnas, detalles y otros.

\begin{figure}[ht]
    \centering
    \includegraphics{images/220px-Drafter_at_work.jpg}
    \caption{Delineante en el trabajo.}
\end{figure}


\item Dibujo mecánico
El dibujo mecánico se emplea en la elaboración de planos para la representación de piezas o partes de máquinas, maquinarias , vehículos como grúas, motos, aviones, helicópteros e industriales. El campo comercial, donde la aplicación practica de los dibujos de ingeniería adopta la forma de dibujos de trabajo, es importante tener en cuenta un amplio conocimiento de los que son los elementos de máquinas, su fabricación y la representación gráfica de cada uno de ellos. Siempre será necesario, que las partes o elementos que ensamblan una maquina se puedan mostrar con facilidad al fabricante y al consumidor, y poder mostrarle con claridad cada una de sus características esenciales y las normas a seguir para la fabricación de cada elemento.

\item Dibujo eléctrico:
Este tipo de dibujo se refiere a la representación gráfica de instalaciones eléctricas en una industria, oficina o vivienda que requiera de electricidad.

\item Dibujo geológico:
El dibujo geológico se emplea en geografía y en geología, en él se representan las diversas capas de la tierra empleando una simbología y da a conocer los minerales contenidos en cada capa. Se usa mucho en minería y en exploraciones de yacimientos petrolíferos.

\item Dibujo topográfico:
Es el dibujo técnico que tiene por objeto representar en un plano las características de un terreno, tales como: el relieve, la altura a diferentes niveles, entre otros. El dibujo topográfico nos representa gráficamente las características de una determinada extensión de terreno, mediante signos convencionalmente establecidos. Nos muestra los accidentes naturales y artificiales, cotas o medidas, curvas horizontales o curvas de niveles.

\item Dibujo urbanístico:
Este tipo de dibujo se emplea en la organización de ciudades: en la ubicación de centros urbanos, zonas industriales, bulevares, calles, avenidas, jardines, autopistas, zonas recreativas entre otros. Se dibujan anteproyectos, proyectos, planos de conjunto, planos de pormenor entre otros.

\item Dibujo técnico de las instalaciones sanitarias:
Tiene por finalidad representar el posicionamiento de cada una de las piezas sanitarias: ducha, lavamanos, retrete, etc. Incluyendo la ubicación de las tuberías internas o externas.

En los planos de instalaciones sanitarias se utilizan diferentes colores normalizados para indicar si se trata de cañerías (tuberías), artefactos o elementos de cloaca, pluvial, ventilación, agua caliente o fría.

También en los cortes se puede observar las pendientes de las cañerías, profundidades, alturas o distancias a planos de referencias.

\item Dibujo técnico electrónico:

Se relaciona con la representación de esquemas y diagramas de circuitos electrónicos de circulación de corriente de poca intensidad tales como radios, televisores, computadoras, entre otros.

\item Dibujo técnico de construcciones metálicas:

Rama del dibujo que se utiliza en el diseño de puentes, galpones, astilleros, herrería en general.

\item Dibujo técnico cartográfico:

Es el empleado en cartografía, topografía y otras representaciones de entornos de gran extensión.
\end{itemize}

\subsubsection*{Formato de papel}

Para plasmar los dibujos en un soporte físico se utilizan formatos de papel de dimensiones normalizadas. Las más utilizadas son la serie A de la norma ISO, cuyos principales tamaños son (medidas en milímetros):
\begin{table}[ht]
    \centering
    \begin{tabular}{|c|c|}\hline
  A0 &	841	x	1189 mm \\
 A1 &	594	x	841 mm \\
 A2 &	420	x	594 mm \\
 A3 &	297	x	420 mm \\
 A4 &	210	x	297 mm \\
 A5 &	148	x	210 mm \\
 A6 &	105	x	148 mm \\
 A7 &	74	x	105 mm \\
 A8 &	52	x	74 mm \\
 A9 &	37	x	52 mm \\
 A10 &	26	x	37 mm \\\hline
    \end{tabular}
    \caption{Formatos de papel ISO}
\end{table}


\begin{figure}[ht]
    \centering
    \includegraphics[width=0.35\textwidth]{images/220px-A_size_illustration2_with_letter_and_legal.svg.png}
    \caption{Tamaños de papel ISO}
\end{figure}
\newpage

\subsubsection*{UNIDAD 1: INTRODUCCIÓN AL DIBUJO TÉCNICO}

El Dibujo Técnico requiere el uso de unos materiales e instrumentos para poder
ejecutarlo luego debemos conocerlos así como sus características y manejo.
De su correcta utilización del compás, la escuadra y el cartabón, la regla y los
diferentes lápices y estilógrafos, depende la calidad, limpieza y perfección del acabado evaluable en esta asignatura.

El mínimo material necesario de dibujo:
\begin{itemize}
    \item Lápiz: en principio se necesitan dos: uno de dureza HB(semiblando) para repasar las soluciones y para la croquización y otro de dureza 2H o
3H(semiduro) para el trazado con instrumental. El uso de portaminas de
diferentes grosores es práctico pero acuérdate de usar las diferentes durezas.
\begin{figure}[ht]
    \centering
    \includegraphics[width=1\textwidth]{images/lapices.jpg}
    \caption{Durezas de lápices}
\end{figure}


\item Compás: mejor que tenga el bloqueo que sirve para abrir o cerrar los brazos y nos da mayor precisión en el trazado y el manejo es más fácil.
\item Regla milimetrada: el tamaño más conveniente es de unos 30 o 50 cm.

\item Escuadra 30 y 45 gratos, mejor no milimetradas y sin escalón, marca
recomendada Faber Castell de unos 21 cm.

\item Papel para dibujo lineal satinado, tamaño DIN-A4
\item Transportador de ángulos y plantilla de curvas es un material opcional
Otros como sacapuntas, mejor metálico y goma blanca para mina blanda y
dura para mina dura.
\end{itemize}


\subsubsection*{Multimedia: Videos}
\begin{center}
\qrcode{https://www.youtube.com/watch?v=9e3BJ9IdELA&list=PLlJ-LmCi75KYjbilDg4P_SPhzdpT4DfG5} \hyperlink{https://www.youtube.com/watch?v=9e3BJ9IdELA&list=PLlJ-LmCi75KYjbilDg4P_SPhzdpT4DfG5}{Videos}

\includegraphics[width=0.5\textwidth]{images/dibujotecnico.png}

\subsubsection{Actividad}
\begin{itemize}
    \item En una construcción grafica cualquiera ¿qué tipo de mina utilizarías, de la serie H
o de la serie B? Explica la respuesta.
\item Los formatos de papel A3 y A4, son los más utilizados en este curso,
especialmente el último. ¿Cuáles son sus medidas?
\item Clasifica los tipos de útiles que se usan en dibujo técnico.
\item Elaborar un resumen sobre que es el dibujo técnico y sus usos y aplicaciones.
\item Copiar el formato de esta hoja A4 en una hoja en blanco, se trata del formato normalizado para la escuela que utilizaremos en dibujo técnico para seguir aprendiendo.
\end{itemize}
\end{center}
\end{document}
